%=======================================================================
\chapter*{Introduction Générale}
\addcontentsline{toc}{chapter}{Introduction Générale}
\setcounter{section}{0}
\renewcommand{\thesection}{\arabic{section}}

% -----------------------------------------------------------------------
\section{Contexte}

L'Imagerie par Résonance Magnétique (IRM) est un outil incontournable pour l'analyse des
pathologies cérébrales. Les tumeurs cérébrales constituent toutefois un défi diagnostique
majeur~: plusieurs types présentent des apparences visuelles proches, et leur interprétation
requiert une expertise médicale approfondie.

L'Intelligence Artificielle offre des perspectives prometteuses pour automatiser cette
analyse. Si le Deep Learning domine la littérature récente, les approches fondées sur des
caractéristiques manuelles (\emph{handcrafted features}) conservent un intérêt fort grâce à
leur interprétabilité, leur faible coût computationnel et leur robustesse sur des jeux de
données de taille modérée.

% -----------------------------------------------------------------------
\section{Problématique}

La classification automatique des tumeurs cérébrales repose sur la qualité des descripteurs
utilisés pour représenter les images. Statistiques d'intensité, textures (GLCM, LBP),
informations fréquentielles (DWT) et contours (HOG) capturent chacun un aspect différent de
l'image~; leur combinaison et leur optimisation conditionnent directement les performances
des modèles de Machine Learning.

Ce projet adopte donc une stratégie d'ingénierie des caractéristiques~: extraire,
optimiser, fusionner et évaluer systématiquement plusieurs familles de descripteurs avant de
comparer les algorithmes de classification.

% -----------------------------------------------------------------------
\section{Objectifs du projet}

L'objectif principal est de développer un système complet de classification multiclasses des
tumeurs cérébrales à partir d'images IRM. Les objectifs spécifiques sont~:
\begin{itemize}[leftmargin=1.4em,itemsep=2pt]
  \item concevoir un pipeline de bout en bout (prétraitement, extraction, classification)~;
  \item optimiser les paramètres des descripteurs GLCM, LBP, DWT et HOG~;
  \item construire et comparer huit ensembles de caractéristiques fusionnés~;
  \item benchmarker huit modèles de Machine Learning et identifier la meilleure configuration~;
  \item valider statistiquement la robustesse des résultats obtenus.
\end{itemize}

% -----------------------------------------------------------------------
\section{Contributions}

Ce travail apporte~:
\begin{itemize}[leftmargin=1.4em,itemsep=2pt]
  \item une optimisation systématique des descripteurs en Phase~1 (score de séparabilité)~;
  \item une fusion multi-branches avec nettoyage des redondances (Phase~2)~;
  \item un benchmark comparatif de 64 expériences (8 modèles $\times$ 8 ensembles)~;
  \item une validation statistique complète (Kappa, Bootstrap, McNemar, Friedman)~;
  \item une analyse d'interprétabilité par importance des branches de caractéristiques.
\end{itemize}

Le Chapitre~1 présente l'état de l'art~; les Chapitres~2 à~5 décrivent le dataset, l'extraction,
l'optimisation et la construction des ensembles~; les Chapitres~6 et~7 rapportent le benchmark
et la validation statistique~; les Chapitres~8 et~9 analysent les résultats et ouvrent vers le
Deep Learning et les modèles Vision-Language.
